\documentclass[10pt]{article}
\usepackage[margin=20mm]{geometry}
\usepackage{amsmath,amssymb,fontspec,unicode-math,xurl,hyperref}
\IfFontExistsTF{DejaVu Serif}{\setmainfont{DejaVu Serif}}{\setmainfont{Latin Modern Roman}}
\setmathfont{latinmodern-math.otf}
\setlength{\parindent}{0pt}\setlength{\parskip}{6pt}\sloppy
\title{Technical review: Proposer-Agent-Evaluator}
\author{AI-assisted technical review}
\date{}
\begin{document}
\maketitle
Abstract: The reported unseen-website evaluation reuses websites from supervised fine-tuning, narrowing a central zero-shot generalization claim. Independent arithmetic and static source checks identify additional local reporting and release defects. Gains relative to the same initialization and the human evaluation evidence remain relevant.\par
\section*{Scope and assessment}
This report reviews I25-43739, Proposer-Agent-Evaluator (PAE): Autonomous Skill Discovery For Foundation Model Internet Agents, by Yifei Zhou and colleagues, in the 39-page PMLR original [paper], together with the acquired December 2024 source release [code]. Assessment: Major for the interpretation of the unseen-website experiment. The paper itself discloses that the supposedly unseen evaluation websites also supplied supervised fine-tuning (SFT) demonstrations. This materially narrows a central zero-shot generalization conclusion. It does not erase the measured improvements over the same SFT initialization.\par
The method uses fixed task-proposer and outcome-evaluator models, a trainable browsing agent, and success-filtered imitation of collected trajectories. There is no proved generalization theorem. The principal dependencies are the training/evaluation split, reward validity, actual optimization objective, and reproducible task coverage. The following findings distinguish manuscript evidence from release-specific implementation behavior.\par
\section*{What the unseen-website experiment establishes}
Problem 1 (Major: the evaluation websites are already represented in SFT).  The abstract and Section 5.3 describe zero-shot generalization to unseen websites; Table 3 calls them websites PAE never interacts with. Appendix F, pages 17 and 22, says the agent is first fine-tuned on all 11,220 Claude-generated trajectories collected from 85 websites outside WebVoyager. Figures 20–21, pages 30–31, explicitly identify those same websites as the websites subsequently used for the generalization evaluation.\par
Thus the task sequence documented by the paper is SFT on these websites, RL on WebVoyager websites, and evaluation back on the SFT websites with 500 synthetic tasks. These domains are unseen during the later RL stage, but not unseen by the trained agent. Domain familiarity can therefore explain part of the resulting performance, and this experiment does not isolate transfer to websites absent from all task-specific training. This is a direct split-definition issue, not a conjecture about pretraining contamination or a claim that identical test instructions were memorized.\par
The strongest alternative reading is that “PAE” denotes only the post-SFT RL stage. Under that definition, Table 3 is a useful test of transfer from newly practiced RL domains back to previously trained domains. The reported \(9.1\%\to16.3\%\) and \(16.1\%\to21.4\%\) changes remain evidence of improvement relative to the SFT baselines. The appendix’s explicit disclosure is also counterevidence to any suggestion of hidden misconduct. It does not support the broader, unqualified claim that the final agents generalize zero-shot to websites never encountered in their training.\par
The repair is either to restrict the conclusion to websites excluded from the RL stage, or to construct a website partition excluded from both SFT and RL and repeat evaluation there. Publish domain, task and trajectory identifiers for each stage so instruction overlap can separately be checked. Figures 20–21 visibly enumerate 13 groups of five URLs, or 65 distinct URLs, rather than the stated 85; the missing list entries also need reconciliation. The independent URL count is saved in the accompanying calculation output.\par
\section*{Numerical and evaluation interpretation checks}
Problem 2 (Minor: some reported gains and weighted averages are inconsistent).  Section 5.3, page 6, reports a 10.8 percentage-point WebArena Easy improvement for the 7B model, but Table 2 gives \(24.6-18.0=6.6\) points. The same paragraph’s 34B WebVoyager gain is 10.4 points, whereas Table 1 gives \(33.0-22.2=10.8\) points. The 10.4-point number correctly describes the separate comparison with Qwen2VL-72B, \(33.0-22.6\).\par
For the stated WebArena Easy counts \((108,50,50)\) and Table 2’s per-domain percentages, the weighted mean must be \[\bar s=\frac{108s_{\rm Map}+50s_{\rm PostMill}+50s_{\rm Market}}{208}.\] The PAE row gives \(25.365\%\), not \(24.6\%\); the user-demo row gives \(26.556\%\), not \(25.7\%\); Claude 3.5 gives \(49.454\%\), not \(50.1\%\). Rounding all component percentages and the final mean to one decimal can create at most 0.1 percentage point of discrepancy. A common running average with fixed domain weights commutes with this weighted sum, so smoothing alone does not resolve the discrepancy. Different valid-task denominators, checkpoints or averaging windows could do so, but must be supplied rather than inferred. The SFT row recomputes to \(18.037\%\), consistent with 18.0. These localized discrepancies do not reverse the displayed SFT-to-PAE improvements.\par
WebVoyager’s Table 1 is explicitly annotated by Claude 3 Sonnet, despite the more general ground-truth wording in Sections 4.1 and 5.1. The independent human study in Figure 4 is valuable counterevidence against dismissing those gains as merely evaluator agreement: on the sampled trajectories the 7B human success rate rises from \(12.5\%\) to \(26\%\). The pooled confusion matrix gives \(91.4\%\) agreement, but positive predictive value is only \(20.6/(20.6+4.2)=83.1\%\) and recall is \(82.4\%\). Small net system bias can coexist with substantial individual mistakes. The displayed seven model pairs give mean absolute system disagreement \(1.93\) points, rather than the prose’s \(1.7\); a different aggregation would need documentation.\par
Needs Evidence: uncertainty and selection. Two hundred human annotations per method, the displayed confusion matrix, and consistent directional gains support the empirical story. They do not establish the significance of a 0.1-point comparison between 7B PAE and 34B SFT, or a 0.7-point comparison against Qwen72B on WebArena Easy. Tables report PAE running averages without a complete smoothing/checkpoint-selection specification or independent training-seed uncertainty. Table 5 lists a hyperparameter search without identifying the selection split. Report raw task-level outcomes, paired comparisons, seed variability, averaging windows and the tuning protocol. No test-set tuning or fabricated result is inferred from missing information. The rewritten 208-task WebArena Easy benchmark must remain distinct from the original benchmark; Appendix G explicitly explains this change and does not report PAE gains on the full original task set.\par
\section*{Optimization: a valid special case and a different released objective}
For binary terminal reward, fresh trajectories sampled from a fixed current policy satisfy the policy-gradient identity \[\nabla J(\theta)=\mathbb E_{\tau\sim\pi_\theta}
 [R(\tau)\nabla\log p_\theta(\tau)]
 =P_{\pi_\theta}(R=1)\,
 \mathbb E_{\pi_\theta(\tau\mid R=1)}[\nabla\log p_\theta(\tau)].\] Thus success-only cloning gives a positive rescaling of the on-policy gradient at the sampling policy, with trajectory/token sums treated consistently and the sampling distribution held fixed for the surrogate gradient. This verifies the intended elementary motivation in Section 4.4. It does not make arbitrary replay-buffer cloning an unbiased current-policy REINFORCE estimator.\par
Problem 3 (Minor, release-scoped: the actual filtering and replay change that objective).  The source appends newly collected trajectories to historical data, retains up to 500,000 trajectories, and clones successful members of that replay buffer without importance weights (onpolicy\_train\_loop.py 171–190). Additionally, utils.py 52–68 retains only the shortest 20th-percentile successes for each task, rather than all successful trajectories as described in Section 4.4.\par
Algorithm 1 itself accumulates a replay buffer, so replay is not wholly undisclosed; it nevertheless requires a qualification of the claimed on-policy equivalence. In a one-step bandit let action \(A\) always succeed and \(B\) succeed with probability \(0.8\). Let the current probability of \(A\) be \(p=0.8\) and a replay-generating policy have \(q=0.1\). All rewards remain binary. With respect to the current logit of \(A\), \[\nabla J=p(1-p)(1-0.8)=0.032,
 \qquad
 \mathbb E_q[R\nabla\log\pi]=q(1-p)-(1-q)0.8p=-0.556.\] Normalizing by replay success probability \(0.82\) preserves the opposite sign. This exact example disproves a general estimator equivalence, not the empirical usefulness of replay. Separately, successful lengths \((2,10)\) have 20th percentile 3.6, so the released filter discards the longer success. Report and ablate this extra length preference, or remove it when claiming literal all-success cloning.\par
The release also averages token log probabilities inside each action (interface.py 133–159), rather than taking the action-sequence log-probability sum in the displayed objective. Gradient clipping with norm 0.01 is listed in the hyperparameter table, but the only clipping call in the inspected BC update is commented out (trainer.py 84–88). The supplied single-server RL configuration uses rollout size 256 and one actor epoch, whereas the table gives 4096 and four for WebVoyager. These may be illustrative release configurations; they are not evidence that the original runs used them. Exact experiment configurations and a distinction between the REINFORCE motivation and the implemented heuristic are needed for faithful reproduction.\par
\section*{Release-specific evaluation defects}
The reviewed commit is `f40715b61993523822174f1b6981fc9a1f2c2dd9`, dated 30 December 2024. Its relationship to the final PMLR experiments has not been established. No author program, remote browser, model or checkpoint was executed.\par
Problem 4 (Minor: the evaluation path does not enumerate every requested task).  scripts/run.py 95 sets the test batch size to \(8\lfloor N/8\rfloor\). BatchedWebEnv.reset enumerates indices from zero, but WebBroswerGym.reset tests `if task\_id`, so index zero falls through to a random task (webgym.py 537–553 and 937–944).\par
For \(N=557\), the path attempts 552 trajectories: indices 1 through 551 plus one uniform random index. Hence at most 552 distinct tasks can be covered, and six indices are reached only through the random slot. The expected distinct count is \(551+6/557=551.0108\). Even \(N=208\), divisible by eight, has full coverage only when the random slot selects zero, with probability \(1/208\); otherwise one task is duplicated and task zero omitted. This follows by enumerating possible random indices in an independent script. It is not a measurement of historical evaluation coverage. Preserve the remainder and test `task\_id is not None`, then verify exact task identifiers in the resulting evaluation log. The inspected main loop keeps evaluation trajectories separate from training trajectories, so this finding does not assert evaluation-data leakage into policy updates.\par
Problem 5 (Minor: verdict parsing uses the entire reasoning text).  utils\_eval.py 566–568 returns success whenever the whole response contains `SUCCESS` and does not contain `NOT SUCCESS`; it does not parse the final verdict requested in Figure 13.\par
The malformed response “The agent was UNSUCCESSFUL.” is accepted because it contains the substring `SUCCESS`. Conversely, reasoning that mentions a hypothetical `NOT SUCCESS` case before concluding `Verdict: SUCCESS` is rejected. Normal canonical outputs are correctly handled, which bounds the issue. Use a structured verdict or an anchored final-field parser and reject malformed outputs explicitly. This identifies a release robustness defect; its occurrence rate and effect on any reported score remain unknown. The evaluator’s last-three-screenshots design is implemented as described, but screenshots and a final answer cannot generally certify hidden historical conditions or every earlier subtask.\par
\section*{Autonomy, visual evidence and surviving conclusions}
The claim that all improvements use only website names needs a stage-specific qualification. Figure 15 explicitly supplies three handwritten OpenStreetMap examples, and Figures 18–19 contain detailed domain instructions. Appendix C says the optional user demonstrations are author-selected pages. Qwen proposer/evaluator variants remove proprietary models from that RL stage, while Appendix F’s Claude-generated SFT remains part of the initialization; they therefore do not demonstrate an entirely proprietary-model-free training pipeline. These disclosed inputs do not negate the absence of human trajectory labels during later autonomous rollouts.\par
Figures 1–7 show the architecture, ablation trends and human checks; Figures 8–12 and 22–30 provide visible examples of navigation, failure and recovery. Selected successful examples are illustrative, not an independent estimate of success rates or a proof of newly acquired causal skills. In particular, Figure 30’s displayed hours menu and final response do not by themselves certify exclusion of all 24-hour businesses. Figure 12 combines a visible HTTP 403 failure with unsupported response content, illustrating that the exclusive error categories can overlap conceptually. Raw annotations and inter-annotator agreement would help interpret the error breakdown; the examples alone cannot validate every category assignment. The reasoning ablation also changes the SFT starting capability in Figure 3, so it does not isolate reflection during RL at matched initial performance.\par
The stronger surviving conclusion is that, in the specified evaluation conditions, an SFT-initialized agent can improve through generated tasks and outcome-filtered practice, including when the proposer/evaluator has lower measured navigation performance than the agent. That empirical observation does not require a new-domain zero-shot interpretation or exact on-policy estimator equivalence. No claim of a universal generalization guarantee, historical code failure, or author correction is made here.\par
Item | Assessment | Repair verification | Author adoption\par
Unseen-domain interpretation | Major | New split needed | Unknown\par
Arithmetic and release details | Minor | Local checks only | Unknown\par
Uncertainty/mechanism scope | Needs Evidence | Raw evidence needed | Unknown\par
\section*{Scope and references}
Main scientific pages 1–9 and the impact statement on page 10 were read. Applicable Appendices A–K on pages 14–39 were read, including the prompts and training-domain lists; image-only trajectories were actually viewed. Relevant source excerpts cover task proposal, train/evaluation separation, replay and filtering, token loss, evaluator parsing, environment task reset and supplied configurations. The dependency installations, all distributed-worker paths, live website state, datasets, model weights, human annotation records, external citations and other manuscript revisions were not exhaustively checked. The separate OpenReview attachment was unavailable in the supplied acquisition record; the embedded appendices were available. There was no empirical replication.\par
The accompanying independent script records its Python/NumPy versions, table inputs, binary-reward counterexample, URL count, index enumeration and parser examples. Source-file hashes and precise excerpts accompany this report. Scientific adjudication remains pending; human reviews are empty and this initial review is not independent acceptance.\par
9 Y. Zhou et al., Proposer-Agent-Evaluator (PAE): Autonomous Skill Discovery For Foundation Model Internet Agents, Proceedings of ICML, PMLR 267 (2025), 79490–79528. Reviewed publisher PDF, https://proceedings.mlr.press/v267/zhou25ah.html. Consulted files are identified above and preserved with hashes.\par
\end{document}
